\section{Related Work}\label{sec:related}

\paragraph{MCP security measurement}
Most large-scale studies capture MCP servers at a single point in time. Three check whether tool descriptions match their implementations. DCIChecker combines static extraction with an LLM judge over 2,214 Python servers and names stale documentation and feature drift as a cause of mismatch, without measuring either~\cite{dcichecker2026}. MCPDiFF compares declared and implemented behavior~\cite{mcpdiff2026}, and MCP-FSD finds file-system capabilities in servers whose category does not imply them~\cite{mcpfsd2026}. Two studies work at ecosystem scale. Parasites in the Toolchain analyzes attacks on the MCP ecosystem~\cite{parasites2026}, and MCPZoo deploys tens of thousands of servers at run time and finds that existing scanners flag 96.89\% of them as risky, with an average alert precision of 45.53\%~\cite{mcpzoo2026}. Two studies add time. Bharti tracks the official registry's short server descriptions over 89 days~\cite{bharti2026registry}, and Kraishan differences registry version records for descriptions, packages, and endpoints~\cite{kraishan2026samename}; both note that tool definitions and their code are not visible in registry data. We measure that surface directly. Our extractor resembles the description-code pipelines, extends them to TypeScript and JavaScript, and adds what none of them measure: how tool-level capabilities change across versions. Because scanner verdicts are this unreliable, we take precision from human labels rather than from an automated judge.

\paragraph{Attacks and defenses}
Tool poisoning hides instructions for the model in tool metadata~\cite{invariant2025poisoning}. Rashidi shows that such payloads reach the model through every metadata surface tested, that invisible Unicode tag characters escape the approval view, and that none of eight tested techniques, including a post-approval rug pull, triggers re-approval in the protocol~\cite{unicodetag2026}. TrustShiftProbe studies servers that turn malicious after gaining trust, using synthetic servers~\cite{trustshiftprobe2026}. Proposed defenses include signed and versioned tool definitions~\cite{etdi2025} and permission manifests for MCP servers~\cite{agentbound2025}. Our data grounds both lines of work in real servers: concealment in descriptions is rare, changes after approval are common, and capability-based checks of the kind these defenses propose would fire on the changes that matter.

\paragraph{Agent skill ecosystems}
Malicious agent skills have been detected in the wild~\cite{liu2026skills}, and concurrent work studies skills over time, tracing how copied skills evolve and gain command execution~\cite{skillconstellations2026} and how revisions to skill files affect agents~\cite{skillevolution2026}. We study MCP servers, whose tools carry code rather than instructions, and measure how that code's capabilities change after a user approves it.

\paragraph{Software supply chain}
Package ecosystem studies show how single maintainers and updates propagate risk~\cite{zimmermann2019npm}, and systematizations catalog attacks through malicious updates and compromised accounts~\cite{ladisa2023sok}. An npm package already runs with the user's authority, so authority is not what MCP adds. What it adds is indirection: a model chooses and invokes tools based on their descriptions without seeing their code, and clients do not ask again when that code changes.

\paragraph{Permissions and warning fatigue}
Mobile permission research found that users rarely read or understand install-time permission prompts~\cite{felt2012android}, and browser warning studies show high click-through when warnings are frequent~\cite{akhawe2013alice}. Our labels show why this matters for MCP: most real changes are routine, so a client that re-asks on every change would mostly interrupt users for nothing. The question is when to re-ask, not only whether.
